\printedpage{24}
% head: A. Grothendieck : Dérivateurs, Ch. XIX, Édition des Universités de Montpellier II et Paris VII
% foot: Version 11 février 2019
\msp{page 43}
instructive, of course, if $\mathcal{M}$ is stable under small $\limind$ , if its objects are accessible etc. and if $W$ is L-accessible (i.e. stable under small filtering $\limind$, and moreover accessible \dots). The first candidate for $I$ that I think of would be a $B_G$ ($G$ a group, e.g. $G=\pm1$), it is indeed the example at the antipodes of ordered sets, which, for their part, do not seem to cause too much trouble (assuming the noetherian hypothesis at least).

The following examples are not all that convincing, alas!

\textbf{Example 3.10.} Let $\mathcal{M}$ be any category, then the couples $(\Fl(\mathcal{M}),\Is(\mathcal{M}))$ and $(\Is(\mathcal{M}),\Fl(\mathcal{M}))$ are both Quillen couples,
\[
\Fl(\mathcal{M}) \Qcouple \Is(\mathcal{M}) \ , \qquad \Is(\mathcal{M}) \Qcouple \Fl(\mathcal{M}) \ . \tag*{(3.10.1)} % reading resolved: DER_XIX_corrections.csv
\]
It follows that if we set
\[
W \ = \ \Is(\mathcal{M}) \ , \qquad C \ = \ F \ = \ \Fl(\mathcal{M}) \ , \tag*{(3.10.2)}
\]
$(W,C,F)$ is a closed Quillen triple \textsuperscript{(}\footnote[22]{at least if we assume that $\mathcal{M}$ is stable under $\limind$ and $\limproj$ that are \textit{finite.}}\textsuperscript{)}. Moreover, for every $I$ in Cat, we have $W(I)=\Is(\mathcal{M}(I))$, hence
\[
\Ho_W(I) \ = \ \mathcal{M}(I) \ . \tag*{(3.10.3)}
\]
So the question of the existence of the functors $u_!^W$, $u_*^W$ reduces to that of the usual functors $u_!$, $u_*$, left and right adjoints to $u^* : \mathcal{M}(I') \longrightarrow \mathcal{M}(I)$.

\msp{page 44}
It is well known that the existence of these functors for a given $u$, calls for the existence of certain $\limind$ resp. $\limproj$ in $\mathcal{M}$; e.g. if $I'=e$, it is necessary and sufficient that the $\limind$ (resp. $\limproj$ ) of type $I^\circ$ be representable in $\mathcal{M}$. But with the hypotheses made on $\mathcal{M}$, there is no problem if $u : I \longrightarrow I'$ is a functor between categories that are \textit{finite.} We are going to give an example where this is no longer so.

\textbf{Example 3.11.} Let $k$ be a commutative ring, and $\mathcal{M}=K^{\mathrm{b}}(k\text{-Ab})$ the category of bounded cochain complexes of $k$-modules, $W$ the set of quasi-isomorphisms, hence
\[
\Ho_W \ = \ D^{\mathrm{b}}(k\text{-Ab}) \ ,
\]
the bounded derived category of the abelian category of $k$-modules. We know, using the existence of enough projectives in $k$-Ab, that this structure is associated with a closed Quillen structure $(W,C,F)$, described by \textsc{Quillen}
\msp{page 45}
in Homotopical Algebra -- except that he takes $D^-(k)$ instead of $D^{\mathrm{b}}(k)$. Perhaps, to verify the axioms, we must assume that $k$ has finite cohomological dimension.

Let us take $I=B_G$, where $G$ is a group (e.g. $G=\{\pm1\}$), and $u : I \longrightarrow e$. Then $\mathcal{M}(I)$ is the category of bounded complexes of $k[G]$-modules. We see that $u_!^W$ resp. $u_*^W$ exists if and only if for every such bounded complex, or simply for a plain $k[G]$-module $M$, the complex $L\Gamma^G(M)$ resp. $R\Gamma_G(M)$ belongs to $D^{\mathrm{b}}(k)$, or what amounts to the same,
\printedpage{25}
% head: A. Grothendieck : Dérivateurs, Ch. XIX, Édition des Universités de Montpellier II et Paris VII
% foot: Version 11 février 2019
if and only if we have
\[
\begin{aligned}
H_i(G,M) \ &= \ 0 \quad \text{for } i \gg 0 \ ,\\
\text{resp.}\quad H^i(G,M) \ &= \ 0 \quad \text{for } i \gg 0 \ ,
\end{aligned}
\]
for every $k[G]$-module $M$. This is practically never the case; when we take e.g. $k=\mathbf{Z}$ and $G$ a finite group, this is the case only if $G=e$. The most economical example seems to me to be $G=\pm1$, $k=\mathbf{F}_2$.

\msp{page 46}
So here $u_!^W$ and $u_*^W$ do not exist, for $I \longrightarrow e$ with $I$ a finite category. But also $\mathcal{M}$ is not stable under small $\limind$ .

\textbf{Example 3.12.} Let $\mathcal{M}$ be a category equipped with a strongly saturated localizer $W$, such that $\mathcal{M} \longrightarrow \underbrace{\mathcal{M}W^{-1}}_{=\mathcal{H}}$ admits an adjoint -- a left one, say. So we have a situation of two adjoint functors % [sic] muni -> munie, agreement with « une catégorie »
\[
\begin{tikzcd}[column sep=huge]
\mathcal{H} \arrow[r, shift left, "v_!\ \text{fully faithful}"] & \mathcal{M} \arrow[l, shift left, "v^*\ \text{localization}"]
\end{tikzcd}
\ , \tag*{(3.12.1)}
\]
and $W=(v^*)^{-1}(\mathrm{Iso}(\mathcal{H}))$. We find that for every $I$ in Cat, we have
\[
\mathrm{HO}_W(I) \ \simeq \ \mathcal{H}(I) \ \overset{\text{def}}{=} \ \mathrm{Hom}(I^\circ,\mathcal{H}). \tag*{(3.12.2)} % [sic] HO_W -> Ho_W, notation Ho_W used in (3.10.3) and Exemple 3.11
\]
So we can think of making an example where in $\mathcal{M}$ the small $\limind$ and $\limproj$ exist, but not in $\mathcal{H}$, so that the $u_!^W$ or $u_*^W$ do not always exist. But this seems difficult to construct, seeing that in the situation envisaged, we easily see that the small $\limproj$ exist in $\mathcal{H}$. So if in $\mathcal{H}$ % [sic] envisagé -> envisagée, agreement with « la situation »
\msp{page 47}
there exists a small cogenerating family (and for this, it is enough that such a family exist in $\mathcal{M}$), there will also be small $\limind$ , by virtue of SGA 4 I 8.12.7. Likewise, they exist if $\mathcal{H}$ is a full subcategory of $\mathcal{M}$ \textit{stable under small} $\limind$ , and it seems practically impossible to make an example (3.12.1) without $\mathcal{H}$ having this stability property. Not to mention that there is even the following perplexity:

\marginpar{! \textbar}\textbf{Question 3.13.} \textit{Is there a $\mathfrak{U}$-category (\,$\mathfrak{U}$ a given universe) in which the small $\limind$ are representable, without the same being true of the small $\limproj$ ?} \textbf{N.B.} Such a category does not admit a small generating family, a fortiori it is not equivalent to a small category.

\msp{page 48}
\textbf{Question 3.14.} \textsc{Anderson} claims that for the prederivator defined by a category with cofibrations satisfying the transfinite axiom Cof$_4$, the functors $u_!^W$ exist. His proof is totally botched, nevertheless I cannot manage to make a counterexample. Could the result he announces be true? I do not even know of an example of a localizer $W$, in a category $\mathcal{M}$ \textit{with small $\limind$ and small $\limproj$ ,} such that $u_!^W$ or $u_*^W$ do not always exist -- but I am convinced that there are examples (even with $\mathcal{M}$ a small category?). I am less convinced of this if we suppose $W$ stable under filtering $\limind$, and if need be, moreover, that the elements of $\mathcal{M}$ are accessible (plus possibly other innocuous categorical properties of $\mathcal{M}$, such as the existence of a cardinal
\printedpage{26}
% head: A. Grothendieck : Dérivateurs, Ch. XIX, Édition des Universités de Montpellier II et Paris VII
% foot: Version 11 février 2019
\msp{page 49}
filtration etc.)

\textbf{Question 3.15.} This brings me back to the particularly elusive and irritating question, of the existence of small Q-generating sets, in a subset $\Phi$ of $\Fl(\mathcal{M})$, when we suppose (say) that $\mathcal{M}$ is stable under small $\limproj$ , that its elements are accessible and that it admits a cardinal filtration. Thus we have the question: does an L-accessible subset of $\Fl(\mathcal{M})$, which we assume if need be to be Q-saturated on the left (or what amounts to the same, containing the isomorphisms, and stable under cobase change, under composition and under direct factors) admit a \textit{small subset} $\Phi_0$ that is Q-generating on the left? % [sic] Q-saturé -> Q-saturée, agreement with « une partie »

\marginpar{! \textbar}On the other hand, I have not at present constructed a single example of a couple $\Phi \longleftrightarrow \Psi$, in a category $\mathcal{M}$, that is not factorizing! The case where $\mathcal{M}$ is special, was of course the most interesting. Clearly, in all
\msp{page 50}
this it is examples that are lacking. The key to an understanding does seem to lie in the questions of factorization between sets $\Phi \longleftrightarrow \Psi$. One ought to make a systematic study of this, but I shrink from the prospect of a digression that would perhaps take me yet more weeks! % [sic] Le clef -> La clef, « clef » is feminine